\scanpage{57}
\diagram{\begin{tikzcd}[ampersand replacement=\&,column sep=small,row sep=small]
 C\arrow[rr,"\varphi"]\&\& E\\
 E\arrow[rr,"{}^a\varphi"]\arrow[dr,dashed,no head]\&\&\widehat C\\
 \&\widehat E\arrow[ur,dashed,"\varphi^*"']\&
\end{tikzcd}}
\[
 \operatorname{Hom}_E(\varphi(Y),X)\simeq\operatorname{Hom}_{\widehat C}(Y,{}^a\varphi(X))
\]
\[
 \begin{aligned}
 &\deleted{\unsure{\operatorname{Hom}_{\widehat E}(\varphi(Y),X)\simeq\operatorname{Hom}_{\widehat C}({}^a\varphi(X),Y)}}\\
 &\deleted{\unsure{\operatorname{Hom}_E(X,}}
 \end{aligned}
\]

Si ds $E$ $\exists$ \inserted{petites} limites inductives, on a aussi :
\[
 \bar\varphi:\widehat C\longrightarrow E\qquad
 \bar\varphi(F)=\varinjlim_{C/F}\varphi(X)
\]
\margin{la donnée de $\varphi$ équivaut à celle de $\bar\varphi$ commutant aux petites $\varinjlim$}
\[
 \operatorname{Hom}_{E}(\bar\varphi(F),X)\simeq\operatorname{Hom}_{\widehat C}(F,{}^a\varphi(X))
\]

\uline{Conditions équivalentes}
\[
 \left.\begin{array}{ll}
 \text{\textcircled{a}}&\forall X\in E,\quad\varinjlim_{Y\in C/X}\varphi(Y)\xrightarrow{\sim}X\\
 \text{\textcircled{b}}&{}^a\varphi:E\longrightarrow\widehat C\ \deleted{\unsure{\mathrm{est}|-}}\ \text{pl. fidèle}
 \end{array}\right\}\quad\text{On dit alors que $\varphi$ est strictement générateur.}
\]
b$'$) (Si $E$ stable par \inserted{petites} $\varinjlim$) $\widehat C\to E$ un foncteur de localisation [ou seulement : ess. surjct

\unsure{c$'$|b$''$|b$'$}) \deleted{\unsure{$\varphi E$}} (Si $E$ est stable par \inserted{petites} $\varinjlim$ filtrantes, et les $C/X$ ($X\in\operatorname{Ob} E$) filtrant, p.ex. $C$ stable par $\varprojlim$ finies et $\varphi$ exact à g :) : $\widetilde\varphi:\operatorname{Ind}(C)\to E$ induit par $\bar\varphi$ est une foncteur de localisation

c) (Si $\varphi$ admet un ind-adjoint, i.e. ${}^a\varphi$ prend ses valeurs dans $\operatorname{Ind}(C)$)
\[
 {}^{\mathrm{ia}}\varphi:E\mathbin{-}\operatorname{Ind}(C)\quad\text{est pl fidèle.}
\]

d) (si $\varphi$ admet un adjoint à droite $\psi$, i.e. ${}^a\varphi$ prend ses valeurs dans $C$)
\[
 \psi\ \text{est pl fidèle}
\]

d$'$) \deleted{$\varphi$ est un fonct\unsure{eur}} (\unsure{m\textsuperscript{e}|même} condition \inserted{que \unsure{f|d}}) $\varphi$ est un foncteur de localisation.

e) (si $\varphi$ pl. fidèle) $\varphi(C)$ est génératrice dans $E$ par épim. stricts.

\textcircled{f} Pour toute catégorie $F$ stable par petites $\varinjlim$, désignant par $\underline{\operatorname{Hom}}'(E,F)$ la sous-catégorie pleine de $\underline{\operatorname{Hom}}$ formée des foncteurs qui commutent aux petites $\varinjlim$, le foncteur restriction
\[
 \beta_F:\underline{\operatorname{Hom}}'(E,F)\mathbin{-}\underline{\operatorname{Hom}}(C,F)
\]
est pl. fidèle

\deleted{g') Avec les notations précédentes, le foncteur $\unsure{\alpha_F}\quad\underline{\operatorname{Hom}}(C,F)$}

g) (si $C$ stable par $\varprojlim$ finies, et $\varphi$ y commute \inserted{et est pl. fid}) le foncteur $\beta_F$ précédent induit une équivalence entre la catégorie $\underline{\operatorname{Hom}}'(E,F)$ et la \deleted{\unsure{cat|sous}} catégorie pleine de $\underline{\operatorname{Hom}}(C,F)$ formée des $f:C\mathbin{-}F$ \deleted{tels} qui commutent à toute \unsure{limite|lim} inductive dans $E$ qui en est une dans $F$

\deleted{\unsure{f)} (si $C$ stable par $\varprojlim$ finies, \inserted{et $\varprojlim$ finies} $C$ admet des petites $\varinjlim$, les \unsure{petites} les petites $\varinjlim$ filtrantes sont exactes, $\varphi$ est exact \unsure{par}}

\unsure{h)} \deleted{\unsure{i.e. que}} (Si $E$ un topos, $C$ stable par $\varprojlim$ finies, $\varphi$ exact à g) $\varphi$ définit $E$ comme \uline{sous-topos} de $\widehat C$.
\deleted{\unsure{g') (même condition que|i) (même condition que}}
% archive sheet number: 56
